% Generated draft: requires mathematical and source review.
\documentclass[11pt]{article}
\usepackage[margin=25mm]{geometry}
\usepackage{fontspec}
\setmainfont{Latin Modern Roman}
\usepackage{amsmath,amssymb,mathtools,unicode-math}
\setmathfont{Latin Modern Math}
\usepackage{enumitem,needspace,xurl,hyperref}
\setlength{\parindent}{0pt}
\setlength{\parskip}{5pt}
\setlength{\emergencystretch}{3em}
\newcommand{\N}{\mathbb{N}}
\newcommand{\Z}{\mathbb{Z}}
\newcommand{\Q}{\mathbb{Q}}
\newcommand{\R}{\mathbb{R}}
\newcommand{\C}{\mathbb{C}}
\newcommand{\F}{\mathbb{F}}
\newcommand{\A}{\mathbb{A}}
\newcommand{\PP}{\mathbb{P}}
\newcommand{\E}{\mathbb{E}}
\newcommand{\eps}{\varepsilon}
\newcommand{\dd}{\,\mathrm{d}}
\begin{document}
\section*{UnsolvedMath ID: 30000885}
\textbf{Source problem code:} catalogue entry
\subsection{Definitions and mathematical statement}
For a finite group $G$, let
\[
\Q G\cong A_1\times\cdots\times A_n
\]
be its Wedderburn decomposition into simple rational algebras. Write $K_i=Z(A_i)$ for the center. An order $R_i\subset A_i$ is a subring with identity that is a full finitely generated $\Z$-lattice in $A_i$. Its reduced-norm-one unit group is
\[
R_i^1=\{u\in R_i^\times:\operatorname{Nrd}_{A_i/K_i}(u)=1\}.
\]
For a virtually torsion-free group $H$, its virtual cohomological dimension is $\operatorname{vcd}(H)=\operatorname{cd}_{\Z}(H_0)$ for a torsion-free subgroup $H_0$ of finite index; it is independent of that choice.

The source reports classifications when
\[
\operatorname{vcd}(R_i^1)\leq2\qquad\text{for every }i.
\]
The requested extension is to the class of finite groups and associated rational components satisfying
\[
\operatorname{vcd}(R_i^1)\leq3\qquad\text{for every }i,
\]
and to obtain a corresponding classification or structural description of
\[
U(\Z G)=(\Z G)^\times.
\]
The condition is componentwise; it is not replaced by the different requirement $\operatorname{vcd}(U(\Z G))\leq3$. The full integral unit group is the target, not only its augmentation-one subgroup.

The extract does not state the existing classification theorems, the intended output (a list of finite groups, allowed simple components, or descriptions up to a specified group equivalence), or whether the unit-group classification is exact or up to finite index/commensurability. No direct-product identity for the integral unit group is inferred merely from the rational Wedderburn product, and the dimension cutoff alone is not presented as a classification theorem.
\subsection{Short English statement}
Extend the known componentwise norm-one-unit classifications from virtual cohomological dimension at most two to at most three, and describe the resulting integral group-ring unit groups. The intended classification equivalence and baseline theorems are not supplied.
\subsection{Sources}
\begin{itemize}
\item Supplied report: \emph{Mini-Workshop: Arithmetik von Gruppenringen} (2008), Oberwolfach Reports, DOI 10.4171/owr/2007/55, pp. 3209--3240; proposed year 2007. \url{https://doi.org/10.4171/owr/2007/55}.
\item Local UnsolvedMath registry, permanent ID 30000885, catalogue entry; created 25 August 2026, updated 16 September 2026. Its corrected\_verified note says the complete extension question and indexed objects were restored. The supplied 21 August 2026 literature triage labels the general extension open, reporting restricted families rather than a full classification. No additional external reference is supplied; this dated assessment is not independently verified here.
\end{itemize}
\end{document}
